\section{Discussion}

\subsection{Key Findings}

\textbf{The EvalPlus failure set is a reproducible foundation for repair research.} The
h-e1-v2 PASS result explicitly verifies reproducibility: failure IDs, stored outputs,
API access, and test selection all confirmed. The 4-condition verification protocol
is a reusable template for establishing reproducibility before executing repair
experiments.

\textbf{Static analysis is the wrong oracle for EvalPlus semantic failures.} The 16--32\%
SA fire rate means two-thirds to five-sixths of failures are opaque to the standard
repair pipeline. Specification-aligned repair --- intent, behavioral gap, deviation ---
is a structured oracle designed for this failure population, not an optional enhancement.

\textbf{The h-e1 $\rightarrow$ h-e1-v2 trajectory illustrates hypothesis refinement.} h-e1 targeted
\texttt{plus\_fail\_tests} completeness (4.5\% gap); h-e1-v2 redesigned to target
\texttt{solutions\_cache.jsonl} coverage (134/134 complete). This was scientifically correct:
the stored incorrect solutions are the relevant data for Condition B/C construction.
Verification failures should prompt target redesign, not hypothesis abandonment.

\subsection{Limitations}

\textbf{Mechanism experiments not executed.} P1, P2, P3 are pre-registered predictions,
not measured results. All comparative claims are INCONCLUSIVE. This paper establishes
infrastructure and design; the comparative results are in the companion mechanism
experiments.

\textbf{6-task exclusion (n$=$128, not n$=$134).} 4.5\% exclusion; principled (tasks cannot
have Condition C prompts from archived data); statistically adequate power maintained.

\textbf{Single model, single temperature.} GPT-4o-mini at temperature$=$0.2, seed$=$42 only.
Generalization to other models, temperatures, and random seeds is future work.

\textbf{Component ablation deferred.} The triple is evaluated as an integrated oracle.
Individual component attribution requires 7 additional conditions; deferred to
future work.

\subsection{Broader Impact}

The 4-condition existence verification protocol advances reproducibility methodology
for LLM code repair evaluation. Pre-registration of P1/P2/P3 before mechanism
experiment execution prevents hypothesis drift. The infrastructure and pre-registration
together constitute an independently valuable scholarly contribution: the verification
protocol is reusable, the archived failure set is publicly documentable, and the
pre-registered design prevents post-hoc analysis bias. No apparent negative societal
impact; improved LLM code repair is broadly beneficial.
